

\begin{figure}[ht]
\begin{center}
\includegraphics[width=0.9\linewidth]{Fig/AME-CAM.png}

\caption{An overview of the proposed AME-CAM method, which contains multiple-exit network based activation extraction phase and attention based activation aggregation phase. The operator $\odot$ and $\otimes$ denote the pixel-wise weighted sum and the pixel-wise multiplication, respectively.
}
\label{Fig.AME-CAM}
\end{center}
\end{figure}

The proposed AME-CAM method consists of two training phases: activation extraction and activation aggregation, as shown in Fig. \ref{Fig.AME-CAM}. In the activation extraction phase, we use a binary classification network, e.g., ResNet-18, to obtain the class probability $y=f(I)$ of the input image $I$. To enable multiple-exit training, we add one internal classifier after each residual block, which generates the activation map $M_i$ of different resolutions. We use a cross-entropy loss to train the multiple-exit classifier, which is defined as
\begin{equation}
\label{eq.mloss}
loss = \sum_{i=1}^4 CE(GAP(M_i), L)
\end{equation}
where $GAP(\cdot)$ is the global-average-pooling operation, $CE(\cdot)$ is the cross-entropy loss, and $L$ is the image-wise ground-truth label.

In the activation aggregation phase, we create an efficient hierarchical aggregation method to generate the aggregated activation map $M_f$ by calculating the pixel-wise weighted sum of the activation maps $M_i$. We use an attention network $A(\cdot)$ to estimate the importance of each pixel from each activation map. The attention network takes in the input image $I$ masked by the activation map and outputs the pixel-wised importance score $S_{xyi}$ of each activation map. We formulate the operation as follows:
\begin{equation}
\label{eq.anet}
S_{xyi} = A([I\otimes n(M_i)]_{i=1}^4)
\end{equation}
where $[\cdot]$ is the concatenate operation, $n(\cdot)$ is the min-max normalization to map the range to [0,1], and $\otimes$ is the pixel-wise multiplication, which is known as image masking. The aggregated activation map $M_f$ is then obtained by the pixel-wise weighted sum of $M_i$, which is $M_f = \sum_{i=1}^4 (S_{xyi} \otimes M_i)$.

We train the attention network with unsupervised contrastive learning, which forces the network to disentangle the foreground and the background of the aggregated activation map $M_f$. We mask the input image by the aggregated activation map $M_f$ and its opposite $(1-M_f)$ to obtain the foreground feature and the background feature, respectively. The loss function is defined as follows:
\begin{equation}
\label{eq.closs}
\begin{aligned}
loss = SimMin(v_i^f, v_j^b)+ SimMax(v_i^f, v_j^f) + SimMax(v_i^b, v_j^b)
\end{aligned}
\end{equation}
where $v_i^f$ and $v_i^b$ denote the foreground and the background feature of the i-th sample, respectively. $SimMin$ and $SimMax$ are the losses that minimize and maximize the similarity between two features (see C$^2$AM\cite{xie2022c2am} for details).

Finally, we average the activation maps $M_1$ to $M_4$ and the aggregated map $M_f$ to obtain the final CAM results for each image. We apply the Dense Conditional Random Field (DenseCRF)\cite{krahenbuhl2011efficient} algorithm to generate the final segmentation mask. It is worth noting that the proposed method is flexible and can be applied to any classification network architecture.

\begin{table}
\centering
\caption{Comparison with weakly supervised methods (WSSS), unsupervised method (UL), and fully supervised methods (FSL) on BraTS dataset with T1, T1-CE, T2, and T2-FLAIR MRI images. Results are reported in the form of mean$\pm$std. We mark the highest score among WSSS methods with bold text.}
\label{Table:Compare_CAM_BraTS}
% \setlength{\tabcolsep}{4pt}
\begin{tabular}{c|c|c|c|l}
\hline
\multicolumn{5}{c}{BraTS T1} \\ \hline
Type           & Method            & Dice $\uparrow$         & IoU $\uparrow$          & \multicolumn{1}{c}{HD95 $\downarrow$}            \\ \hline
\multirow{6}{*}{WSSS}           & Grad-CAM (2016)   & 0.107$\pm$0.090 & 0.059$\pm$0.055 & 121.816$\pm$22.963 \\
               & ScoreCAM (2020)   & 0.296$\pm$0.128 & 0.181$\pm$0.089 & 60.302$\pm$14.110  \\
               & LFI-CAM (2021)    & 0.568$\pm$0.167 & 0.414$\pm$0.152 & 23.939$\pm$25.609  \\
               & LayerCAM (2021)   & 0.571$\pm$0.170 & 0.419$\pm$0.161 & 23.335$\pm$27.369  \\
               & Swin-MIL (2022)   & 0.477$\pm$0.170 & 0.330$\pm$0.147 & 46.468$\pm$30.408  \\
               & AME-CAM  (ours)   & \textbf{0.631$\pm$0.119} & \textbf{0.471$\pm$0.119} & \textbf{21.813$\pm$18.219}  \\ \hline
UL             & C\&F (2020)       & 0.200$\pm$0.082 & 0.113$\pm$0.051 & 79.187$\pm$14.304  \\ \hline
\multirow{2}{*}{FSL}            & C\&F (2020)       & 0.572$\pm$0.196 & 0.426$\pm$0.187 & 29.027$\pm$20.881  \\
               & Opt. U-net (2021) & 0.836$\pm$0.062 & 0.723$\pm$0.090 & 11.730$\pm$10.345  \\ \hline
\multicolumn{5}{c}{BraTS T1-CE} \\ \hline
Type           & Method            & Dice $\uparrow$         & IoU $\uparrow$          & \multicolumn{1}{c}{HD95 $\downarrow$}            \\ \hline
\multirow{6}{*}{WSSS}           & Grad-CAM (2016)   & 0.127$\pm$0.088 & 0.071$\pm$0.054 & 129.890$\pm$27.854 \\
               & ScoreCAM (2020)   & 0.397$\pm$0.189 & 0.267$\pm$0.163 & 46.834$\pm$22.093  \\
               & LFI-CAM (2021)    & 0.121$\pm$0.120 & 0.069$\pm$0.076 & 136.246$\pm$38.619 \\
               & LayerCAM (2021)   & 0.510$\pm$0.209 & 0.367$\pm$0.180 & 29.850$\pm$45.877  \\
               & Swin-MIL (2022)   & 0.460$\pm$0.169 & 0.314$\pm$0.140 & 46.996$\pm$22.821  \\
               & AME-CAM  (ours)   & \textbf{0.695$\pm$0.095} & \textbf{0.540$\pm$0.108} & \textbf{18.129$\pm$12.335}  \\ \hline
UL             & C\&F (2020)       & 0.179$\pm$0.080 & 0.101$\pm$0.050 & 77.982$\pm$14.042  \\ \hline
\multirow{2}{*}{FSL}            & C\&F (2020)       & 0.246$\pm$0.104 & 0.144$\pm$0.070 & 130.616$\pm$9.879  \\
               & Opt. U-net (2021) & 0.845$\pm$0.058 & 0.736$\pm$0.085 & 11.593$\pm$11.120  \\ \hline
\multicolumn{5}{c}{BraTS T2} \\ \hline
Type           & Method            & Dice $\uparrow$         & IoU $\uparrow$          & \multicolumn{1}{c}{HD95 $\downarrow$}            \\ \hline
\multirow{6}{*}{WSSS}           & Grad-CAM (2016)   & 0.049$\pm$0.058 & 0.026$\pm$0.034 & 141.025$\pm$23.107 \\
               & ScoreCAM (2020)   & 0.530$\pm$0.184 & 0.382$\pm$0.174 & 28.611$\pm$11.596  \\
               & LFI-CAM (2021)    & 0.673$\pm$0.173 & 0.531$\pm$0.186 & 18.165$\pm$10.475  \\
               & LayerCAM (2021)   & 0.624$\pm$0.178 & 0.476$\pm$0.173 & 23.978$\pm$44.323  \\
               & Swin-MIL (2022)   & 0.437$\pm$0.149 & 0.290$\pm$0.117 & 38.006$\pm$30.000  \\
               & AME-CAM  (ours)   & \textbf{0.721$\pm$0.086} & \textbf{0.571$\pm$0.101} & \textbf{14.940$\pm$8.736}   \\ \hline
UL             & C\&F (2020)       & 0.230$\pm$0.089 & 0.133$\pm$0.058 & 76.256$\pm$13.192  \\ \hline
\multirow{2}{*}{FSL}            & C\&F (2020)       & 0.611$\pm$0.221 & 0.474$\pm$0.217 & 109.817$\pm$27.735 \\
               & Opt. U-net (2021) & 0.884$\pm$0.064 & 0.798$\pm$0.098 & 8.349$\pm$9.125    \\ \hline
\multicolumn{5}{c}{BraTS T2-FLAIR} \\ \hline
Type           & Method            & Dice $\uparrow$         & IoU $\uparrow$          & \multicolumn{1}{c}{HD95 $\downarrow$}            \\ \hline
\multirow{6}{*}{WSSS}           & Grad-CAM (2016)   & 0.150$\pm$0.077 & 0.083$\pm$0.050 & 110.031$\pm$23.307 \\
               & ScoreCAM (2020)   & 0.432$\pm$0.209 & 0.299$\pm$0.178 & 39.385$\pm$17.182  \\
               & LFI-CAM (2021)    & 0.161$\pm$0.192 & 0.102$\pm$0.140 & 125.749$\pm$45.582 \\
               & LayerCAM (2021)   & 0.652$\pm$0.206 & 0.515$\pm$0.210 & 22.055$\pm$33.959  \\
               & Swin-MIL (2022)   & 0.272$\pm$0.115 & 0.163$\pm$0.079 & 41.870$\pm$19.231  \\
               & AME-CAM  (ours)   & \textbf{0.862$\pm$0.088} & \textbf{0.767$\pm$0.122} & \textbf{8.664$\pm$6.440}    \\ \hline
UL             & C\&F (2020)       & 0.306$\pm$0.190 & 0.199$\pm$0.167 & 75.651$\pm$14.214  \\ \hline
\multirow{2}{*}{FSL}            & C\&F (2020)       & 0.578$\pm$0.137 & 0.419$\pm$0.130 & 138.138$\pm$14.283 \\
               & Opt. U-net (2021) & 0.914$\pm$0.058 & 0.847$\pm$0.093 & 8.093$\pm$11.879   \\ \hline
\end{tabular}
\end{table}

